\section{Analyse d'interprétabilité}
\label{sec:analysis}

\subsection{Importance des caractéristiques}

Pour comprendre quelles familles de descripteurs pilotent les prédictions de Full(opt),
nous avons calculé des scores d'importance à partir d'une forêt aléatoire entraînée
sur le jeu complet, avec regroupement des caractéristiques par branche source.
La Figure~\ref{fig:importance} montre la contribution relative de chaque famille.

\begin{figure}[H]
  \centering
  \includegraphics[width=0.80\linewidth]{rf_branch_importance.png}
  \caption{Importance des caractéristiques par branche (forêt aléatoire sur Full(opt)).
           Le HOG et le LBP contribuent le plus ; les statistiques et la DWT apportent
           des gains complémentaires plus modestes.}
  \label{fig:importance}
\end{figure}

Le HOG domine en raison de sa haute dimensionnalité et de sa forte discriminabilité
basée sur les gradients. Le LBP apporte une information texturale complémentaire.
Les statistiques et la DWT fournissent des gains plus faibles mais non nuls, en
cohérence avec l'amélioration observée du HOG seul (90,3\,\%) au Full(opt) (92,0\,\%).

\subsection{Visualisation de l'espace des caractéristiques}

La Figure~\ref{fig:tsne} présente une projection t-SNE 2D des vecteurs Full(opt)
à 526 dimensions pour tous les échantillons de test.

\begin{figure}[H]
  \centering
  \includegraphics[width=0.70\linewidth]{tsne_best_feature_set.png}
  \caption{Visualisation t-SNE des caractéristiques Full(opt) (partition de test).
           Une séparation partielle des classes est visible ; les clusters gliome
           et méningiome se chevauchent davantage que l'hypophysaire, ce qui explique
           la confusion résiduelle du meilleur modèle.}
  \label{fig:tsne}
\end{figure}

La tumeur hypophysaire forme un cluster bien séparé, ce qui explique la forte
précision par classe pour cette catégorie. Le chevauchement partiel entre gliome
et méningiome correspond à la principale zone de confusion dans les résultats
de classification.
